% CIKM 2026 Demonstration draft, rewritten 2026-09-30 around the layered
% retrieval ladder (catalog → descriptions → wide recall) + fail-closed
% engine-graded judgment; related-work section added; all measured numbers
% re-sourced to the fresh monitored run via the \NUM* macros (PROVISIONAL
% until scripts/107 + audit fill them). Earlier pass: 2026-09-23 re-source to
% run experiment_chat_20260922-211000.
% Official format: sigconf, non-anonymous, <=4 pages of body incl. appendices/acks.
% Only GenAI Usage Disclosure and references may extend beyond that limit.
% AUTHOR ACTION: replace author placeholders; confirm submission eligibility.
\documentclass[sigconf]{acmart}
\setcopyright{none}
\settopmatter{printacmref=true,printfolios=true}
\acmConference[CIKM '26]{Proceedings of the 35th ACM International Conference on Information and Knowledge Management}{November 9--11, 2026}{Rome, Italy}
\acmBooktitle{CIKM '26, November 9--11, 2026, Rome, Italy}
\acmISBN{}
\acmDOI{}

\usepackage{booktabs}
\usepackage{xurl}
\newcommand{\sys}{\textsc{Qdcvr}}
% ── Measured-run number macros (2026-09-30 rewrite) ─────────────────────────
% Every \NUM* macro is filled from the monitored rerun
% benchmark-suite/results/runs/run-20260930T150453Z-3ec2174 via
% scripts/107_exp_grade.py (grounded/keyword-verified/latency/cost) +
% audit_paper_numbers.py + the citation regex over saved answers
% (9/6/1). Values are script-asserted, not hand-typed; audit with
% grep 'NUM' and re-run the graders after any rerun.
\newcommand{\NUMgroundedA}{10}\newcommand{\NUMgroundedB}{10}\newcommand{\NUMgroundedC}{8}
\newcommand{\NUMkwA}{10}\newcommand{\NUMkwB}{8}\newcommand{\NUMkwC}{6}
\newcommand{\NUMciteA}{9}\newcommand{\NUMciteB}{6}\newcommand{\NUMciteC}{1}
\newcommand{\NUMlatA}{116\,s}\newcommand{\NUMlatB}{48\,s}\newcommand{\NUMlatC}{41\,s}
\newcommand{\NUMcostA}{\$0.55}\newcommand{\NUMcostB}{\$0.63}\newcommand{\NUMcostC}{\$0.28}
% One full-width figure per page; Figs. 1/2/3 are double-column floats placed
% by the dbltop queue. Target placement pp. 2/3/4 so the body (through the
% Conclusion) ends on page 4 and only the disclosure + references follow.
\setcounter{topnumber}{1}
\setcounter{totalnumber}{4}
\renewcommand{\topfraction}{.9}
\renewcommand{\dbltopfraction}{.9}
\setcounter{dbltopnumber}{1}
\renewcommand{\textfraction}{.05}
\renewcommand{\floatpagefraction}{.8}

\begin{CCSXML}
<ccs2012>
<concept><concept_id>10002951.10003260</concept_id><concept_desc>Information systems~Information retrieval</concept_desc><concept_significance>500</concept_significance></concept>
<concept><concept_id>10003120.10003121</concept_id><concept_desc>Human-centered computing~Interactive systems and tools</concept_desc><concept_significance>300</concept_significance></concept>
</ccs2012>
\end{CCSXML}
\ccsdesc[500]{Information systems~Information retrieval}
\ccsdesc[300]{Human-centered computing~Interactive systems and tools}
\keywords{knowledge-base management, retrieval-augmented generation, evidence inspection, agent workflows, demonstration}

\begin{document}
\title{QDCVR: Content-Routed Organization and Layered Agent Retrieval with Inspectable Evidence}
% CIKM Demo is single-blind. These are explicit draft placeholders, not anonymity.
\author{Author Name}
\affiliation{\institution{Affiliation}\city{City}\country{Country}}
\email{author@example.com}
\renewcommand{\shortauthors}{Author Name}
% sec0_abstract.tex — 2026-09-30 rewrite around the layered ladder +
% engine-graded judgment storyline. Numbers flow from the \NUM* macros in
% main.tex (filled from the monitored rerun; see PROVISIONAL note there).
\begin{abstract}
\emergencystretch 1.5em\relax
Research administrators must organize shared paper collections and inspect
the evidence behind the answers drawn from them. We present \sys{}, an
agent-operated platform that organizes by content and retrieves layer by
layer through one shared address space. Packaged agent skills parse
documents, route them by content into category bases, and index the parts
under base\slash document\slash part\slash section addresses. Retrieval walks a layered
ladder --- catalog, shelf descriptions, wide vector recall --- and
acceptance is decided per segment by a locally hosted judgment engine
that reads candidates and grades them against the query: the engine is
fail-closed, similarity never overrules its verdicts, and every segment
graded positive survives into the answer. When no rung finds a plausible
candidate, the search stops there and returns a scoped not-found report.
The platform's own ingest loop built a 100-paper, five-base instance;
ingest, storage, and retrieval probes then audited it. In a monitored
three-mode run under one harness, model, and neutral prompt, the platform
grounds \NUMgroundedA{} of 10 designated-paper answers and
keyword-verifies \NUMkwA{} of 10 against fixed gold key facts, with
part/section-level citations in \NUMciteA; the bare agent grounds
\NUMgroundedB{} but verifies \NUMkwB, and dense retrieval grounds
\NUMgroundedC{} and verifies \NUMkwC{} with \NUMciteC. A 27-run
three-lane stability suite shows behaviorally consistent retrieval and
honest abstention whenever the corpus does not cover a question. Per
question the platform costs \$0.55, on par with the bare agent's
file-tool loop (\$0.64) and about twice dense retrieval (\$0.28), while
being the only mode whose evidence is inspectable and addressable.
Attendees organize a
document, watch the ladder retrieve, and examine a failure report. Demo
video:
\url{https://github.com/kingdol666/rag-knowledge/blob/master/paper_demo/video/qdcvr-demo.mp4}.
\end{abstract}

\maketitle

% Full-width floats are declared up front so they meet pages 2/3/4 in reading
% order; declaring them later lets them drift past the experiment section.
\begin{figure*}[t]
\centering
\includegraphics[width=0.80\textwidth]{../figures/submission-20260922/fig1-core.pdf}
\caption{QDCVR at a glance, as two four-panel storyboards. Top row: a
conventional dense-RAG configuration --- documents are shredded into fixed
chunks, two excerpts are retrieved and packed into one context, and a single
generation produces an answer whose sources are not addressable. Bottom
row: the two QDCVR agent loops --- autonomous ingest (parse, route by
content, tag, index parts with base/document/part/section addresses) and
layered retrieval (recall through catalog, descriptions, and a wide
vector net; the judgment engine grades every candidate segment; only
survivors are read, yielding a cited answer or a scoped not-found
report). The top row is for contrast.}
\Description{Two storyboard rows of four numbered panels each. Top row
(baseline): documents arrive in the shared pool; they are shredded into
fixed 800-character chunks, with torn-evidence marks; two excerpts are
retrieved and packed into one 4,000-character context; one generation
returns an answer with no sources attached. Bottom row (QDCVR): a PDF is
parsed and routed by content into the right base; every part is indexed
under the layered address base, document, part, section; the agent walks
the ladder --- catalog, shelf descriptions, wide vector recall --- and the
judgment engine grades each candidate's segments, keeping those it grades
positive; the outcome is a cited answer with a section reference or a
scoped not-found report. A takeaway strip states that the same base,
document, part, and section identifiers power both loops, so every answer
can be checked at its source.}
\label{fig:teaser}
\end{figure*}

\begin{figure*}[t]
\centering
\includegraphics[width=0.84\textwidth]{../figures/submission-20260922/fig2-architecture.pdf}
\caption{Platform operations and the agent-executed QDCVR workflow, with
live-console screenshots (1 collection; 2 search; 3 per-channel scores,
base/part address, verification status; 4 the agent's answer beside the
judgment layer's outcomes). Bottom: the ORGANIZE and RETRIEVE loops.}
\Description{Two areas. The top area shows four live-console screenshots
from the demonstration instance: a knowledge-base row with its document
count; the search form with strategy chips and the NISQ query; a result's
combined and vector scores with its base chip, part address, and a
content-verification status; and the agent's answer after retrieval,
opening a numbered list of the Gene Ontology's three namespaces. Beside
the answer sit the judgment layer's outcomes: segments graded positive
survive and the answer cites them; candidates with nothing positive are
discarded; when nothing survives, one librarian pass through catalog and
descriptions backs up the vector lane, and still nothing surviving yields
a scoped not-found report; dashed outlines mark the fallback exits. The
bottom area shows the two packaged loops: organize runs from a prepared
PDF pool through parsing and a 30,000-character split, routes by content,
and tags and indexes parts; retrieve rewrites the query, recalls across
catalog, descriptions, and vectors, grades candidate segments with the
judgment engine, and reads only the survivors. A note records that the
same 42 kb\_* tools run over vector, graph, and file storage.}
\label{fig:arch}
\end{figure*}
% Fig. 3 declared here (right after Fig. 2) so the dblfloat queue
% meets it in time for a page-4 top placement.
\begin{figure*}[t]
\centering
\includegraphics[width=0.82\textwidth]{../figures/submission-20260920/fig3-answers-dual.pdf}
\caption{Two questions from the 100-paper corpus under the monitored
three-mode run (one harness, model, and neutral prompt; one run per
mode; ellipses elide text): each column shows the saved timeline, a
verbatim excerpt, and sources. Panel~1 (BQ04): all three modes answer
--- A quotes the source at part~1 of~2, section~2; B names the file; C
quotes the chunk copy. Panel~2 (BQ06): A has the judgment engine grade
candidate segments and returns a keyword-verified answer quoting the
source; C reaches the same factor from chunk recall alone, with no
grading step and no addresses.}
\Description{Two stacked panels. Panel one, the Gene Ontology question,
has three columns. Column A's timeline lists one subagent probe, one
vector search, one cross-check search, and two document reads; its
excerpt names Molecular Function, Biological Process, and Cellular
Component and quotes the primer; its sources name the life-sciences base,
the primer paper, part 1 of 2, section 2, and a verbatim quote. Column
B's timeline lists a glob, two greps, and a read; its excerpt states the
same three ontologies; its source names the paper file. Column C's
timeline lists a harness call and one vector search; its excerpt states
the three ontologies; its source names the chunk copy in the
800-character chunk index. Panel two, the precipitation-extremes
question, has two columns. Column A's timeline lists four vector-search
recalls, a shelf check, one judgment-engine call, and a harness call; its
excerpt quotes the saturation-vapor-pressure passage; its sources
name the climate base, the paper, engine-graded reading, and keyword
verification. Column C's timeline lists a subagent probe and five vector
searches; its excerpt names the thermodynamic contribution with a
precision caveat; its source shows the chunk index with no grading step.}
\label{fig:threeway}
\end{figure*}
% sec1_intro.tex — 2026-09-30 rewrite: layered retrieval (逐级检索) as the
% central contribution; three design decisions re-anchored to the current
% agent contracts (layered ladder + engine-graded judgment + honest failure).
\section{Introduction}
\label{sec:intro}

Maintaining a shared paper collection is organizational work before it is
retrieval work: someone must keep the corpus ordered and be able to show
which passages support an answer. Packed-context RAG collapses retrieval
into an opaque context: provenance is lost at pack time, similarity stands
in for reading, and a failed retrieval passes
silently~\cite{rag-lewis,dpr}. Recent systems repair parts of this
pipeline --- Self-RAG and Corrective RAG add self-assessment and
correction~\cite{selfrag,crag}, DeepRead grounds answers in document
hierarchy~\cite{deepread}, and RAPTOR builds hierarchical
indexes~\cite{raptor} --- but each treats the corpus as given. To our
knowledge, none manages the corpus that retrieval consumes through the
same addresses that answers cite.

\sys{} (Query-Driven Content-Verified Retrieval) is an agent-operated
platform that closes this gap with three design decisions. \emph{One
address space}: parsed documents are indexed as
base\slash document\slash part\slash section, and the same identifiers serve organization,
search, and citation. \emph{The ladder proposes; the judge disposes}:
recall runs through a layered ladder --- catalog, shelf descriptions, wide
vector nets --- but acceptance is decided per segment by a locally hosted
judgment engine that reads the candidate's actual content and grades it
against the query; the engine's verdicts are fail-closed, similarity never
overrules them, and every segment graded positive is kept.
\emph{Failure is a first-class output}: when a ladder rung finds no
plausible candidate the search stops there and reports honestly ---
shelves scanned, descriptions read, nothing matched --- instead of
answering from a nearest neighbor.

Both halves are executed by packaged agent skills over one tool surface:
the organization skill parses PDFs, routes them by content into category
bases, and tags and indexes the parts; the retrieval skills walk the
ladder, apply the graded judgment, and cite the passages that survive it
(Figures~\ref{fig:teaser} and~\ref{fig:arch}). The rest of the paper
details the loops, walks the demonstration, and reads a monitored
three-mode run as operational evidence: on ten fixed questions, only the
platform both grounds and keyword-verifies \NUMgroundedA{} of 10 answers,
and a 27-run stability suite shows the three retrieval lanes staying
behaviorally consistent and honestly abstaining whenever the corpus does
not cover a question.

% sec1_related.tex — 2026-09-30 new section (user request: explicit related
% work in the storyline). Compact by design; the page budget belongs to the
% system and evaluation sections.
\section{Related Work}
\label{sec:related}

\textbf{Packed-context RAG.} Dense retrievers feed top-$k$
chunks to the generator~\cite{rag-lewis}; provenance is lost at pack time
and a failed recall is invisible downstream. \sys{} keeps retrieval and
reading as separate, recorded steps.

\textbf{Self-assessing and corrective RAG.} Self-RAG trains reflection
tokens~\cite{selfrag}; Corrective RAG routes by a retrieved-document
evaluator~\cite{crag}. Both judge \emph{retrieval self-reports};
\sys{}'s judgment layer reads candidate content with a locally hosted
engine whose unavailability stops the pipeline rather than degrading it.

\textbf{Hierarchical and graph structures.} RAPTOR recurses summaries
into a tree index~\cite{raptor}; DeepRead navigates document
hierarchy~\cite{deepread}; GraphRAG derives community
graphs~\cite{graphrag}. These structures serve retrieval, but
the corpus remains given; in \sys{} the hierarchy is created by the
organization loop and the same base\slash document\slash part\slash
section identifiers carry ingest placement, retrieval reads, and answer
citations.

\textbf{Abstention.} Calibration and selective answering are studied at
the generator~\cite{kamath}. \sys{} makes abstention structural: a rung
with no plausible candidate is a terminal, reportable state, measured as
such in Section~\ref{sec:eval}.



% sec2_system.tex — 2026-09-30 rewrite: the retrieval ladder (逐级检索) is the
% core subsection. Contract facts cross-checked against the current packaged
% skills (knowledgebase-search / knowledgebase-librarian / knowledgebase-hybrid)
% and the live MCP inventory (42 kb_* tools in a 95-tool inventory).
\section{System and Agent Workflow}
\label{sec:system}

\subsection{Services and Interfaces}
\sys{} binds Section~\ref{sec:intro}'s decisions to one platform: a web
console, a FastAPI backend, a \texttt{ragctl} CLI, and an MCP
server~\cite{mcp} exposing 42 \texttt{kb\_*} tools inside a 95-tool
inventory. Console, CLI, and MCP clients operate the same resources;
only the packaged skills add policy. ChromaDB stores BGE-M3 vectors~\cite{bge-m3}, Neo4j
holds optional document relations, and Markdown plus JSON/YAML metadata
keep the collection inspectable.

\subsection{Organization: From PDF to Addressable Parts}
MinerU converts PDFs to structured Markdown~\cite{mineru}; texts
beyond a 30,000-character ceiling split into parts with line-offset
reads. The organization skill inspects the parsed
content, selects among existing bases, proposing a new one only when
none fits, and places and tags the document. Each placement ends in a final check (three of nine
recorded checks), so a misroute is flagged at ingest, not at answer
time. Administrators keep
control of bases; five bases is a demonstration configuration, not a
limit.

% Pipeline audit table — declared early so the column float lands on p3,
% keeping page 4 for Figure 3 plus the evaluation and conclusion text.
\begin{table}[t]
\caption{Pipeline audit on the 100-paper corpus. Every value is
script-asserted from saved artifacts, not hand-counted; the content audit
samples the first 1{,}000 characters per part.}
\label{tab:pipeline}
\small
\setlength{\tabcolsep}{4pt}
\begin{tabular}{@{}lc@{}}
\toprule
\textbf{Stage / check} & \textbf{Result} \\
\midrule
Papers parsed, routed by content & 100/100, 5 bases \\
Stored parts after 30k split ceiling & 322 \\
Per-paper ingest checks (nine recorded checks) & 100/100 pass \\
True duplicates & 0 (35 near-pairs) \\
Descriptions upgraded & 165/165 \\
Content + storage audits & 322/322 \,\textbullet\, 33/33 \\
Retrieval regression probes & 15/15 \\
Dense replica index (800-char chunks) & 4{,}528 \\
\bottomrule
\end{tabular}
\end{table}

\subsection{Retrieval: The Layered Ladder}
\label{sec:ladder}
The retrieval skills run one layered spine (Figure~\ref{fig:arch},
bottom), and any harness may delegate it to a single retrieval subagent
that returns document identifiers and evidence notes while the main agent
answers.
\emph{L0 catalog.} One call reads every base's name, description, and
document count; a base is pruned only when its description is present,
judged trustworthy, and clearly unable to cover the query.
\emph{L1 descriptions.} One call per candidate shelf collects document
identifiers with their descriptions; a keyword pass over titles and
descriptions supplements them. Descriptions are treated as claims, not
facts --- they narrow the candidates but never decide acceptance.
\emph{L2 query shaping.} The query is rewritten as a declarative
sentence with keywords; literal enumeration questions take a machine-side
literal scan, because semantic recall under-recalls exact mentions.
\emph{L3 judgment.} Candidates are passed to the judgment layer \emph{by
reference}: a locally hosted engine fetches each document, segments it,
and scores every segment against the rewritten query. The engine runs
fail-closed --- if it is unavailable the pipeline stops, it does not
degrade to similarity-only acceptance --- and every segment graded
positive survives. A
0.95-similarity hit whose segments all fail is discarded; a 0.40 hit with
one positive segment survives.
\emph{L4 reading.} Only survivors are read at their recorded addresses
(base, document, part, section), with key numbers quoted verbatim.
\emph{L5 answer.} The answer cites what it read; the record keeps the
search path, backend and threshold, confidence, and blind spots.

Early exit is part of the spine: if L0--L1 surface no plausible candidate
--- or a literal scan hits nothing on a kept shelf --- the search stops
and returns a scoped not-found report; minutes of doomed re-searching are
a contract violation, not thoroughness.

\subsection{Lanes: Fast Recall, Full Recall, Fusion}
The ladder is the contract; three lanes trade breadth against latency.
The \emph{fast lane} rewrites the query, casts a wide vector net (high
$k$, balanced across bases, snippets only), cross-checks the top
candidates with a second, BM25-style pass~\cite{bm25}, and feeds the
deduplicated candidates to the same judgment layer --- ranking proposes,
the engine still disposes; when cross-check reads confirm a candidate
at its recorded address, the lane answers directly. The \emph{librarian lane} is the ladder above
with no vector calls, for unknown shelves or complete-recall questions.
The \emph{fusion lane} runs both concurrently and merges on normalized
(base, document) identity. Judgment is by reference, so document text
stays server-side and the context receives one merged evidence pack.
\emph{Content-verified} denotes this engine judgment, not guaranteed
correctness; saved artifacts keep excerpts, verdicts, and answer
records.


% sec3_demo.tex — 2026-09-30 rewrite: scenarios re-anchored to the ladder
% (organize / inspect a layered retrieval / examine a failure report).
% Trace-sourced numbers marked NUM* are filled from the monitored rerun.
\section{Interactive Demonstration}
\label{sec:demo}

The demonstration targets research administrators and agent developers.
A prepared instance, built by the
ingest loop of Section~\ref{sec:system}, holds 100 papers in five bases
(ingest audit: Table~\ref{tab:pipeline}); attendees drive an agent
with the packaged skills over the kb\_* tools. A live
answer takes from about twenty seconds to about two minutes, with a
six-minute tail on the hardest question; presenters replay saved traces
alongside live queries, the video walks the scenarios, and the repository
README reproduces the instance.

\paragraph{1. Organize, then 2.\ inspect an answer's evidence.}
Scenario one places a PDF: the trace records the parse, proposed base,
tags, and indexing checks.
Scenario two asks which three ontologies subdivide the Gene Ontology. In
the saved trace (BQ04; Section~\ref{sec:eval}), the fast lane recalls
the primer with a vector search plus a BM25-style cross-check, reads the
part at its line offset, and the agent answers with a verbatim quotation
whose source block names base, document, part, and section (``2.~What is
the Gene Ontology?''). The cross-check reads confirmed this candidate directly; the graded
path appears in Figure~\ref{fig:threeway}, panel~2. The cited part opens in the
console at its line offset.

\paragraph{3. Examine a failure report.}
Next, the attendee asks for training epochs in the fictitious paper
\emph{Spectral Tuning for Low-Resource Odor Recognition}. In about
thirty seconds the skill returns a conclusion-first not-found report
backed by its searched-facts list: two-stage recall (ten generic hits,
none the paper), vector search (top similarity 0.52), metadata search
(twelve weak title-token matches), and a judgment-engine call with
nothing to adjudicate. BERT fine-tuning recipes appear among the
near-misses and are not substituted; the remedy offered is re-search by
exact identifier, and nothing is speculated.




% sec4_eval.tex — 2026-09-30 rewrite; compacted 2026-10-01.
% (1) monitored three-mode rerun on the current contract (run
%     benchmark-suite/results/runs/run-20260930T150453Z-3ec2174; graded by
%     scripts/107_exp_grade.py, audited by experiments/audit_paper_numbers.py);
% (2) the 27-run stability suite (review/stability-multiround-20260930).
\section{Evaluation}
\label{sec:eval}

\textbf{Setup.}
The platform's own ingest loop built the corpus, so records exist for
every stage (Table~\ref{tab:pipeline}). Ten fixed questions
(BQ01--BQ10), each answerable only from its
designated paper, ran through three modes differing only in evidence
route, permissions, and prompt framing (one harness, model, and neutral
prompt): the platform answers via its fast lane over the kb\_* tools
(Section~\ref{sec:ladder}); the bare agent is the same model with file
tools over the exported papers; dense retrieval answers through a
single vector-search tool over the 4{,}528-chunk replica
(Table~\ref{tab:pipeline}).
Metrics: \emph{grounded} (the trace cites the designated paper, strict
arXiv-id); \emph{keyword-verified} ($\ge$60\% of fixed gold key
facts at word boundaries); \emph{citations} (part/section/table
markers in saved answers).

\textbf{Findings.}
The platform grounds \NUMgroundedA/10 designated-paper answers and
keyword-verifies \NUMkwA/10 --- a clean sweep under conservative grading;
the bare agent grounds \NUMgroundedB/10 yet verifies \NUMkwB/10; dense
retrieval grounds \NUMgroundedC/10 and verifies \NUMkwC/10 (both dense
grounding misses were tool-availability aborts). Citations naming
part, section, or table reach
\NUMciteA/10 versus \NUMciteB{} (bare) and \NUMciteC{} (dense)
(Figure~\ref{fig:threeway}). The bare agent always locates the designated
paper, yet its verification trails and its answers name files where the
platform names checkable addresses: locating, synthesizing, and citing
are different skills --- the split is this demo's design difference.

\textbf{Honest failure and stability.}
In a 27-run stability suite (found / not-covered / out-of-corpus tasks
--- attention, seismic-wave classification, pickle production ---
$\times$ fast, ladder, and fusion lanes $\times$ three rounds) on the
live multi-base instance, no run fabricated: the found task produced the
correct attention formula in 9/9 runs; the not-covering task answered 9/9
that the library holds no such content; the out-of-corpus task refused
9/9, the librarian lane stopping at the catalog layer. Session-trace
replay confirms 27/27 lane-contract compliance, with median wall-clock
34--76\,s across cells
--- and three scripted out-of-corpus questions against the audited
instance all receive scoped not-found reports (Section~\ref{sec:demo}
shows one).

\textbf{Costs and limitations.}
The platform is slowest (\NUMlatA{} per question versus
\NUMlatB/{}\NUMlatC{} bare/dense) --- the gap buys multi-step reading
and graded judgment --- but not costliest: \NUMcostA{} per question,
about the bare agent's \NUMcostB{} and twice dense retrieval's
\NUMcostC. Grading is conservative: word-boundary matching scores the
bare agent's plural ``cathodes'' (BQ08) at zero. Two platform cells
re-ran once after a harness timeout. Correctness calls are the authors'
trace readings; cache state was uncontrolled; the suite's $\pm$35\,s
variance sits in the answer-generation window.

% sec5_conclusion.tex — 2026-09-30 rewrite (ladder + judge storyline);
% compacted 2026-10-01 for the four-page body budget.
\section{Conclusion}
\label{sec:conclusion}

\sys{} carries organization and retrieval through one address space,
with acceptance as a fail-closed engine judgment. The monitored run
yields grounded, keyword-verified answers in ten of ten cases; the
lanes stay consistent across 27 runs; abstention behaves as designed;
every number re-derives from saved
artifacts.

% Any author-supplied acknowledgments must fit within the four-page body limit.
\section*{GenAI Usage Disclosure}
OpenAI Codex and GLM (via ZCode) assisted with prose, figures, artifact
inspection, and the monitored reruns of Section~\ref{sec:eval}; those
runs' saved traces are retained artifacts, and the analysis and manuscript
remain human-checked. The demonstrated system uses language models for
content routing, query rewriting, and answer drafting; segment-level
acceptance is decided by a locally hosted judgment engine. These outputs
are not independent correctness assessments. Responsibility for the final
claims, sources, figures, and manuscript rests with the human authors.
Competing interests: none declared. Demo video:
\url{https://github.com/kingdol666/rag-knowledge/blob/master/paper_demo/video/qdcvr-demo.mp4}.
Code, traces, and scripts:
\url{https://github.com/kingdol666/rag-knowledge}.

\bibliographystyle{ACM-Reference-Format}
\bibliography{refs}
\end{document}






